% LaTeX 仿真报告 - 12-bit 50MS/s SAR ADC PVT & Noise Analysis
% 编译方式: xelatex -> xelatex (两次)
\documentclass[journal]{IEEEtran}

\usepackage{amsmath,amssymb}
\usepackage{graphicx}
\usepackage{booktabs}
\usepackage{siunitx}
\usepackage{textcomp}
\usepackage{array}
\usepackage{xeCJK}
\setCJKmainfont{SimSun}
\setCJKmonofont{SimSun}

\hyphenation{op-tical net-works semi-conduc-tor}

\sisetup{
  detect-all = true,
  per-mode = symbol,
  number-unit-product = \,,
}

\graphicspath{{./figures/}}

\begin{document}

\title{12-bit 50MS/s SAR ADC CAAZ \\ PVT Corner \& 噪声分析仿真报告}

\author{
  赵~毅,~\IEEEmembership{Huazhong University of Science and Technology}
}

\markboth{仿真实验记录}{12-bit SAR ADC PVT \& Noise Analysis Report}

\maketitle

\begin{abstract}
本报告记录了 12-bit 50MS/s SAR ADC（CAAZ 架构）在 TSMC 18RF 工艺下的 PVT 角仿真结果。
共 5 个工艺角（TT/FF/SS/SF/FS），包含瞬态噪声仿真（100$\times$ 噪声功率）及无噪声对照实验。
报告涵盖 SNR、ENOB、关键模块功耗及噪声分析，并对比不同采样时间（\emph{td}=26.5\,ns vs 24.05\,ns）
下的噪声性能差异。
\end{abstract}

\begin{IEEEkeywords}
SAR ADC, CAAZ, PVT corner, transient noise, $kT/C$ noise cancellation, ENOB.
\end{IEEEkeywords}

\section{仿真条件}
\label{sec:conditions}

\subsection{电路与工艺}
\begin{itemize}
  \item \textbf{电路:} 12-bit 50MS/s SAR ADC with CAAZ (Current-Adjusted AutoZeroing)
  \item \textbf{工艺:} TSMC 18RF (1P6M, 1.8V)
  \item \textbf{架构特性:} Split-sampling CDAC, $kT/C$ noise cancellation,
  \item \hphantom{\textbf{架构特性:}} 单位电容 $C_F = \SI{5}{fF}$，CDAC 总电容 $\approx \SI{100}{fF}$
\end{itemize}

\subsection{仿真参数}
\begin{table}[!t]
\centering
\caption{仿真条件汇总}
\label{tab:conditions}
\renewcommand{\arraystretch}{1.2}
\begin{tabular}{ll}
\toprule
\textbf{参数} & \textbf{设定值} \\
\midrule
仿真器 & Spectre 231 (APS) \\
分析类型 & transient + noise \\
仿真时长 & \SI{15}{\us} \\
时钟频率 & \SI{50}{\MHz} \\
输入信号 & 差分正弦，$f_{in}=61/128\times 10\,\text{MHz}=\SI{4.765625}{\MHz}$ \\
输入幅度 & $v_a = \SI{890}{mV}$，共模 $V_{CM} = \SI{900}{mV}$ \\
采样时间 & $t_d = \SI{26.5}{ns}$（默认）；$\SI{24.05}{ns}$（对比） \\
\bottomrule
\end{tabular}
\end{table}

\subsection{瞬态噪声设置}
\begin{table}[!t]
\centering
\caption{瞬态噪声参数}
\label{tab:noise_settings}
\renewcommand{\arraystretch}{1.2}
\begin{tabular}{ll}
\toprule
\textbf{参数} & \textbf{设定值} \\
\midrule
noisefmax & \SI{10}{GHz} \\
noiseseed & 64261112 \\
noisescale & 10（$100\times$ 噪声功率） \\
噪声使能器件 & I1（时钟反相器链）, I14（CMOS 采样开关） \\
errpreset & moderate \\
cmin & \SI{1.23}{aF} \\
\bottomrule
\end{tabular}
\end{table}

\subsection{工艺角定义}
\begin{table}[!t]
\centering
\caption{5 个 PVT 角定义}
\label{tab:pvt_defs}
\renewcommand{\arraystretch}{1.2}
\begin{tabular}{clll}
\toprule
\textbf{Point} & \textbf{Corner} & \textbf{Temp} & \textbf{Sections} \\
\midrule
1 & TT & \SI{27}{\celsius} & tt + tt\_res + tt\_mim \\
2 & FF & \SI{-40}{\celsius} & ff + ff\_res + ff\_mim \\
3 & SS & \SI{150}{\celsius} & ss + ss\_res + ss\_mim \\
4 & SF & \SI{27}{\celsius} & sf + sf\_res + sf\_mim \\
5 & FS & \SI{27}{\celsius} & fs + fs\_res + fs\_mim \\
\bottomrule
\end{tabular}
\end{table}

\section{PVT 角仿真结果}
\label{sec:pvt_results}

\subsection{仿真概况}
\begin{table}[!t]
\centering
\caption{各 Point 完成状态}
\label{tab:completion}
\renewcommand{\arraystretch}{1.2}
\begin{tabular}{clrrr}
\toprule
\textbf{Point} & \textbf{Corner} & \textbf{Data Size} & \textbf{Elapsed} & \textbf{Errors/Warnings} \\
\midrule
1 & TT & 672\,MB & 15\,min & 0 err, 6 warn, 16 notices \\
2 & FF & 691\,MB & 16\,min & 0 err, 1 warn, 9 notices \\
3 & SS & 660\,MB & 15\,min & 0 err, 1 warn, 9 notices \\
4 & SF & 649\,MB & 15\,min & 0 err, 1 warn, 9 notices \\
5 & FS & 664\,MB & 15\,min & 0 err, 1 warn, 9 notices \\
\bottomrule
\end{tabular}
\end{table}

\subsection{SNR \& ENOB}
\label{subsec:snr_enob}

\begin{table}[!t]
\centering
\caption{各工艺角 SNR/ENOB（$t_d=\SI{26.5}{ns}$, 100$\times$ 噪声）}
\label{tab:snr_enob}
\renewcommand{\arraystretch}{1.2}
\begin{tabular}{lSSS}
\toprule
\textbf{Corner} & \textbf{SNR (dB)} & \textbf{SNDR (dB)} & \textbf{ENOB (bit)} \\
\midrule
TT & \multicolumn{3}{c}{\footnotesize 需在 Virtuoso ADE/MAESTRO GUI 中查看} \\
FF & \multicolumn{3}{c}{\footnotesize 同上} \\
SS & \multicolumn{3}{c}{\footnotesize 同上} \\
SF & \multicolumn{3}{c}{\footnotesize 同上} \\
FS & \multicolumn{3}{c}{\footnotesize 同上} \\
\bottomrule
\end{tabular}
\end{table}

% SNDR_DAC waveform fitting info
% 注：SNDR_DAC 使用 B<11:0> 数字码重构模拟输出，与实际采样值对比计算

\subsection{功耗分析}
\label{subsec:power}

\begin{table}[!t]
\centering
\caption{各工艺角关键模块功耗（\si{\micro\watt}）}
\label{tab:power}
\renewcommand{\arraystretch}{1.2}
\begin{tabular}{lSSS}
\toprule
\textbf{Corner} & \textbf{Comparator} & \textbf{Switch Array} & \textbf{SR Latch} \\
\midrule
TT & \multicolumn{3}{c}{\footnotesize 需在 Virtuoso ADE/MAESTRO GUI 中查看} \\
FF & \multicolumn{3}{c}{\footnotesize 同上} \\
SS & \multicolumn{3}{c}{\footnotesize 同上} \\
SF & \multicolumn{3}{c}{\footnotesize 同上} \\
FS & \multicolumn{3}{c}{\footnotesize 同上} \\
\bottomrule
\multicolumn{4}{l}{\footnotesize 注：Preamp 功耗因 Maestro 表达式设置问题，数据异常，未列入。}
\end{tabular}
\end{table}

\subsection{Preamp 功耗说明}
预放大器（Preamp）的功耗测量在 Maestro 表达式中存在设置问题（ADE-3022 警告：
instance name mismatch between schematic and outputs setup），导致测量结果异常。
该数据已被排除，待修正网表或表达式后补充。

\section{无噪声对照仿真}
\label{sec:no_noise}

\subsection{实验设计}
在 TT corner 下，移除瞬态噪声参数（noisefmax, noiseseed, noisescale, noiseon），
重新运行瞬态仿真，以评估噪声对性能的影响。

\begin{table}[!t]
\centering
\caption{无噪声 vs 带噪声（TT, $t_d=\SI{26.5}{ns}$）}
\label{tab:noise_comparison}
\renewcommand{\arraystretch}{1.2}
\begin{tabular}{lSS}
\toprule
\textbf{Metric} & \textbf{With Noise} & \textbf{Without Noise} \\
\midrule
SNR (dB) & \multicolumn{2}{c}{\footnotesize 需在 Maestro GUI 中查看} \\
SNDR (dB) & \multicolumn{2}{c}{\footnotesize 同上} \\
ENOB (bit) & \multicolumn{2}{c}{\footnotesize 同上} \\
Comparator Power (\si{\micro W}) & \multicolumn{2}{c}{\footnotesize 同上} \\
Switch Array Power (\si{\micro W}) & \multicolumn{2}{c}{\footnotesize 同上} \\
SR Latch Power (\si{\micro W}) & \multicolumn{2}{c}{\footnotesize 同上} \\
\bottomrule
\end{tabular}
\end{table}

\begin{table}[!t]
\centering
\caption{无噪声仿真数据量}
\label{tab:nonoise_size}
\renewcommand{\arraystretch}{1.2}
\begin{tabular}{lrr}
\toprule
\textbf{仿真} & \textbf{数据量} & \textbf{耗时} \\
\midrule
td=26.5ns 带噪声（参考） & 672\,MB & 15\,min \\
td=26.5ns 无噪声 & 418\,MB & 9\,min \\
\bottomrule
\end{tabular}
\end{table}

\section{采样时间对比分析}
\label{sec:td_comparison}

\subsection{实验设计}
将采样时间 $t_d$ 从 \SI{26.5}{ns} 调整为 \SI{24.05}{ns}（TT corner, 100$\times$ 噪声），
评估 CAAZ 电路在不同采样时间下的噪声抑制效果。

\begin{table}[!t]
\centering
\caption{$t_d=\SI{24.05}{ns}$ 仿真结果（TT, 100$\times$ 噪声）}
\label{tab:td2405}
\renewcommand{\arraystretch}{1.2}
\begin{tabular}{lSS}
\toprule
\textbf{Metric} & \textbf{$t_d=\SI{26.5}{ns}$} & \textbf{$t_d=\SI{24.05}{ns}$} \\
\midrule
SNR (dB) & \multicolumn{2}{c}{\footnotesize 需在 Maestro GUI 中查看} \\
SNDR (dB) & \multicolumn{2}{c}{\footnotesize 同上} \\
ENOB (bit) & \multicolumn{2}{c}{\footnotesize 同上} \\
Comparator Power (\si{\micro W}) & \multicolumn{2}{c}{\footnotesize 同上} \\
Switch Array Power (\si{\micro W}) & \multicolumn{2}{c}{\footnotesize 同上} \\
SR Latch Power (\si{\micro W}) & \multicolumn{2}{c}{\footnotesize 同上} \\
\bottomrule
\end{tabular}
\end{table}

\begin{table}[!t]
\centering
\caption{td=24.05ns 仿真数据量}
\label{tab:td2405_size}
\renewcommand{\arraystretch}{1.2}
\begin{tabular}{lrr}
\toprule
\textbf{仿真} & \textbf{数据量} & \textbf{耗时} \\
\midrule
td=24.05ns 带噪声 & 547\,MB & 12\,min \\
td=24.05ns 无噪声 & 412\,MB & 9\,min \\
\bottomrule
\end{tabular}
\end{table}

\subsection{噪声功率削减分析}
基于 $t_d=\SI{26.5}{ns}$ 和 $t_d=\SI{24.05}{ns}$ 的 SNR 差异，计算噪声功率相对变化：

\begin{equation}
\Delta \text{SNR} = \text{SNR}_{24.05} - \text{SNR}_{26.5} \quad [\text{dB}]
\label{eq:delta_snr}
\end{equation}

\begin{equation}
\text{Noise Power Ratio} = 10^{-\Delta \text{SNR} / 10}
\label{eq:noise_ratio}
\end{equation}

噪声功率削减因子 = $\displaystyle\frac{P_{n,26.5}}{P_{n,24.05}}$

\section{综合分析}
\label{sec:analysis}

\subsection{PVT 稳健性}
5 个工艺角全部顺利完成（0 错误），各角数据量约 650-690\,MB，仿真耗时约 15 分钟。
无噪声仿真（td=26.5n）数据量 418\,MB 相比带噪声（672\,MB）减少约 38\%，
因无需计算噪声贡献，仿真速度提升约 40\%。

\subsection{噪声抑制效果}
CAAZ 电路仅对采样开关（I14）和时钟链（I1）引入的 $kT/C$ 噪声有抑制效果。
关闭噪声后，系统热噪声基底下降，SNR/ENOB 提升预计...
（具体数值需在 Maestro GUI 中查看表达式结果）

\subsection{$t_d$ 对噪声抑制的影响}
更长的采样时间提供更充分的噪声消除，但也会影响 SAR 逻辑的可用转换时间。
td=24.05ns vs 26.5ns 的噪声对比需在 Maestro 中查看 SNR 结果后计算。

\section{结论}
\label{sec:conclusion}
已完成 TSMC 18RF 工艺下 12-bit 50MS/s SAR ADC 的 PVT 仿真（5 角）、
无噪声对照仿真及采样时间对比仿真。全部 8 组仿真均成功完成（0 errors），
总数据量 4.7\,GB。具体 SNR/ENOB/功耗数据需在 Virtuoso ADE/MAESTRO GUI 中查看表达式计算结果。

\end{document}
